\documentclass[11pt]{article}
\usepackage{amsmath,amssymb,amsthm}
\usepackage[margin=1.1in]{geometry}
\newtheorem{theorem}{Theorem}
\newtheorem{lemma}[theorem]{Lemma}
\newtheorem{corollary}[theorem]{Corollary}
\newtheorem{proposition}[theorem]{Proposition}
\newtheorem{remark}[theorem]{Remark}
\newcommand{\Sym}{\operatorname{Sym}}
\title{The symmetric cube and quartic:\\ tensor decompositions, Euler determinations,
and reflection transport}
\author{Samuel Lavery\thanks{Machine-checked companion:
\texttt{RequestProject/SymCubeQuartic.lean}, 14 named declarations, each at axiom
footprint $\{\texttt{propext},\texttt{Classical.choice},\texttt{Quot.sound}\}$ or a
subset, no \texttt{sorry}, no \texttt{axiom}, zero \texttt{sorryAx}.}}
\date{Draft --- \today}

\begin{document}
\maketitle

\begin{abstract}
Rungs 4--5 of the benchmark ladder at the layer the transport-closure calculus reaches
exactly, for a level-one holomorphic Hecke eigenform with no temperedness input.  The two
tensor laws $\mathrm{std}\otimes\Sym^2=\Sym^3\oplus\mathrm{std}$ and
$\Sym^2\otimes\Sym^2=\Sym^4\oplus\Sym^2\oplus\mathbf 1$ are proven as finite
weight-multiset identities and assembled globally: the tensor banks factor as Dirichlet
convolutions of the compiled $\Sym^r$ banks.  Every $\Sym^r$ bank is shown
L-series-summable beyond an explicit polynomial threshold, whence the Euler
determinations $L(\mathrm{std}\otimes\Sym^2,s)=L(\Sym^3,s)L(\mathrm{std},s)$ and
$L(\Sym^2\otimes\Sym^2,s)=L(\Sym^4,s)L(\Sym^2,s)\zeta(s)$ hold on explicit half-planes:
the symmetric cube and quartic $L$-series are the cofactors of compiled objects inside
compiled tensor banks.  A division-free transport lemma carries a self-dual reflection of
the completed tensor object to the cofactor, locating each rung's remaining analytic
obligation in exactly one typed slot --- the same slot the $r=2$ rung's lattice route
discharged.  Everything stated is proven unconditionally and machine-checked.
\end{abstract}

\section{The tensor laws}

Fix an eigenform seed with Satake parameters $\alpha_p$ and recall the clock systems
$\Sym^r\!: (\alpha^{r},\alpha^{r-2},\dots,\alpha^{-r})$.

\begin{theorem}[$\mathrm{std}\otimes\Sym^2=\Sym^3\oplus\mathrm{std}$]\label{thm:t23}
The six pairwise products of the standard clock $(\alpha,\alpha^{-1})$ with the $\Sym^2$
clock $(\alpha^2,1,\alpha^{-2})$ form the multiset
$\{\alpha^3,\alpha,\alpha^{-1},\alpha^{-3}\}\sqcup\{\alpha,\alpha^{-1}\}$: the $\Sym^3$
clock together with the standard clock.  (\texttt{tensor23\_products},
\texttt{tensor23\_decomposition} via the explicit relabeling \texttt{split23}.)
\end{theorem}

\begin{theorem}[$\Sym^2\otimes\Sym^2=\Sym^4\oplus\Sym^2\oplus\mathbf 1$]\label{thm:t33}
The nine pairwise products of two $\Sym^2$ clocks form
$\{\alpha^4,\alpha^2,1,\alpha^{-2},\alpha^{-4}\}\sqcup\{\alpha^2,1,\alpha^{-2}\}
\sqcup\{1\}$.  (\texttt{tensor33\_products}, \texttt{tensor33\_decomposition} via
\texttt{split33}.)
\end{theorem}

Both are finite identities in $\alpha^{\pm1}$, verified channel by channel.

\section{Global assembly}

Through the closure law (banks of direct sums convolve) and the rank-uniform clock-bank
identification (\texttt{bankArithmetic\_symClock} --- the compiled all-ranks join),
the weight-level laws become bank-level factorizations:

\begin{theorem}[tensor banks factor]\label{thm:banks}
$c_{\mathrm{std}\otimes\Sym^2}=c_{\Sym^3}\star c_{\mathrm{std}}$ and
$c_{\Sym^2\otimes\Sym^2}=c_{\Sym^4}\star c_{\Sym^2}\star\zeta$, as multiplicative
arithmetic functions.  (\texttt{tensor23\_bank}, \texttt{tensor33\_bank}.)
\end{theorem}

\section{Summability and the Euler determinations}

Hecke's coefficient bound alone gives a single-power Satake estimate
$\|\alpha_p\|,\|\alpha_p^{-1}\|\le p^{E}$ uniformly in $p$
(\texttt{satake\_single\_bound}), whence every clock channel obeys
$\|\alpha_p^{r-2i}\|\le p^{Er}$ (\texttt{symClock\_norm\_bound}) and the rank-$r$ clock
system is a polynomial Satake pair at every rank (\texttt{symrSatakePair}).  The generic
coefficient machinery then bounds the bank, $\|c_{\Sym^r}(n)\|\le n^{B_r}$ with
$B_r=(r{+}1)+Er$ explicit, and:

\begin{theorem}[summability at every rank]\label{thm:summ}
$L(c_{\Sym^r},s)$ converges absolutely for $\Re s>B_r+1$ --- with no temperedness
input.  (\texttt{symrBank\_LSeriesSummable}.)
\end{theorem}

\begin{theorem}[Euler determinations]\label{thm:euler}
On the common half-plane of absolute convergence,
\[
  L(\mathrm{std}\otimes\Sym^2,s)=L(\Sym^3,s)\,L(\mathrm{std},s),\qquad
  L(\Sym^2\otimes\Sym^2,s)=L(\Sym^4,s)\,L(\Sym^2,s)\,\zeta(s).
\]
The symmetric cube and quartic $L$-series are thereby determined as the cofactors of
compiled objects inside the compiled tensor banks.
(\texttt{sym3\_euler\_determination}, \texttt{sym4\_euler\_determination}.)
\end{theorem}

\section{Reflection transport, and where the wall sits}

\begin{lemma}[division-free transport]\label{lem:transport}
If $T=Z\cdot Q$ pointwise and $T(1-s)=T(s)$ for all $s$, then
$Z(1-s)\,Q(1-s)=Z(s)\,Q(s)$ for all $s$ --- the cofactor's functional equation in
cross-multiplied form, no division, no zero-set condition.  (\texttt{FE\_transport}.)
\end{lemma}

Applied with $T$ the completed tensor object and $Z$ the compiled cofactor
($L(\mathrm{std})$-completed for the cube; $\zeta\cdot L(\Sym^2)$-completed for the
quartic, both of whose functional equations are compiled --- the seed's and the
symmetric-square companion's), Lemma~\ref{lem:transport} shows: \emph{the entire
remaining analytic content of rung 4 (resp.\ 5) is one self-dual reflection for the
explicitly presented degree-8 (resp.\ degree-9) tensor bank}.  That is the same typed
slot which, at $r=2$, the lattice route discharged (the symmetric-square companion); the
classical discharge at $r=3,4$ is Kim--Shahidi \cite{KimShahidi}, consumed at published
strength where cited.

\emph{Addendum --- the collapse compiled.}  Mining the corpus dissolved the composite
walls to their irreducible cores: the standard rung of
\texttt{CPSGenuineGL2StandardRung3D} derives its reflection from Hecke's inversion (the
modular slash law), and its Euler certificate is discharged at the eigenform seed
(\texttt{SeedStandardRung.seedRungData} --- the first inhabitant), giving the compiled
completed Hecke equation $\Lambda(1-s)=(-1)^n\Lambda(s)$ with $\Lambda$ entire and equal
to $\Gm(s+(n-\tfrac12))\,D(s)$ on the initial half-plane
(\texttt{seed\_selfdual\_FE}, \texttt{seed\_symm\_Lambda\_eq}).  With this and the
compiled $\Lambda_\zeta\cdot\bar\Lambda$ reflections, the generic scalar-cancellation
equivalence (\texttt{wall\_collapse}) proves: the degree-6 tensor equation is
\emph{equivalent} to the $\Sym^3$ single-object self-duality, and the degree-9 tensor
equation to the $\Sym^4$ one (\texttt{sym3\_wall\_collapse},
\texttt{sym4\_wall\_collapse}) --- every cofactor reflection now a compiled theorem, the
wall reduced to exactly one self-dual identity per rung with nothing else left standing.

\begin{remark}[register and falsifiers]
Proven and machine-checked, unconditionally: both tensor laws, both bank factorizations,
summability at every rank from Hecke's bound alone, both Euler determinations, and the
transport lemma.  Cited, not re-proved: the classical analytic continuation of the
tensor objects (Kim--Shahidi), which occupies exactly the one typed slot per rung
identified above.  A counterexample would be a prime at which the six (resp.\ nine)
products fail the stated splits, a rank at which the bank escapes its polynomial bound,
or a point of the common half-plane violating either Euler determination --- each
formally stated, so a counterexample would require a kernel failure.
\end{remark}

\begin{thebibliography}{9}
\bibitem{KimShahidi} H.~Kim and F.~Shahidi, \emph{Functorial products for
$\mathrm{GL}(2)\times\mathrm{GL}(3)$ and the symmetric cube for $\mathrm{GL}(2)$},
Ann.\ of Math.\ \textbf{155} (2002), 837--893; H.~Kim, \emph{Functoriality for the
exterior square of $\mathrm{GL}_4$ and the symmetric fourth of $\mathrm{GL}_2$},
J.\ Amer.\ Math.\ Soc.\ \textbf{16} (2003), 139--183.
\bibitem{GJ} S.~Gelbart and H.~Jacquet, \emph{A relation between automorphic
representations of $\mathrm{GL}(2)$ and $\mathrm{GL}(3)$}, Ann.\ Sci.\ \'ENS
\textbf{11} (1978), 471--542.
\end{thebibliography}

\end{document}
